\begin{papercard}{grundy2019}
\cf{Stages} \stagetag{A} \stagetag{D} \stagetag{H} \stagetag{R}. Traffic, content and network analysis of 24 prominent Android medicines apps. It traces user data to first, third and ``fourth'' parties and models aggregation and re-identification potential across the ecosystem.

\cf{Aim} Investigate whether and how top-rated medicines apps share user data and characterise privacy risks for clinicians and consumers \pc{grundy2019}{1}{To investigate whether and how user data are shared}. Two phases: (1) traffic analysis; (2) content and network analysis of recipients \pc{grundy2019}{2}{We carried out this study in two phases}.

\cf{Sampling (purposive, triangulated)} ``Prominent'' meant highly downloaded, top-100 ranked, or endorsed by credible organisations. From 17 October to 17 November 2017 two strategies were triangulated: (a) a crawler sampling the weekly top-100 free and paid apps of the Medical category in the UK, US, Australian and Canadian Google Play stores; (b) recommended or endorsed apps from a not-for-profit medicines organisation, a curated app library, a systematic review and pharmacists' networks \pc{grundy2019}{2}{We purposefully sampled medicines related apps that}. One investigator screened 821 app names for medicines relevance. Two investigators independently screened 67 store descriptions against the criteria: medicines-related, Android/Australia availability, at least one ``dangerous'' permission or a claim to collect data, and interactivity requiring user input \pc{grundy2019}{3}{One investigator screened 821 apps for any app}. Exclusions: single-company or single-country restriction, prohibitive cost, unavailable. The flow diagram shows 24 included apps \pc{grundy2019}{3}{Sampling flow diagram for prominent medicines related apps}.

\cf{Traffic protocol} Agrigento (detects obfuscated and encrypted leaks) on a Pixel 1 with Android 7.1. Subscriptions were purchased if needed. Four scripted dummy profiles (one doctor, one pharmacist, two consumers) were randomised to apps; all buttons were clicked, all settings adjusted and profile data entered \pc{grundy2019}{3}{we simulated real world, in-depth use of the app using}. Differential analysis: 14 baseline runs, then a 15th run with one of 28 prespecified data types altered (e.g.\ Android ID, birthday, email, location). A deviation counts as a leak \pc{grundy2019}{3}{Using one randomly assigned dummy scripted}. Recipients were identified with Shodan and WHOIS.

\cf{Content analysis of recipients} For each recipient, two investigators independently examined the Crunchbase profile, website, policies, terms and investor prospectus, extracting mission, activities, data-sharing partnerships and privacy practices into a REDCap form \pc{grundy2019}{3}{For each of the entities receiving user data in the traffic}. Extraction happened before and after the GDPR came into force, and some policies disclosed more partnerships afterwards. Discrepancies were resolved by consensus, taking the newer information \pc{grundy2019}{4}{Data were extracted between 1 February}. Entities were coded inductively in NVivo by main activity and business model.

\cf{Party typology and network} First parties (developer/parent), users as second parties, third parties (direct external recipients), \emph{fourth parties} (companies with which third parties report further sharing) \pc{grundy2019}{4}{We classified entities receiving user data into three}. Two networks in R/igraph: direct recipients, and potential ecosystem sharing including fourth parties. A simulation assumed one person uses all apps, to show how data multiply and which firms can aggregate \pc{grundy2019}{4}{This simulation assumes that the same person uses all}. Patient and public involvement came through a consumer advocacy network partner.

\cf{Key findings} 19/24 apps shared user data. 55 entities from 46 parent companies received data; 67\,\% were analytics/advertising-type services. Third parties advertised sharing with 216 fourth parties, and some firms held central aggregating positions \pc{grundy2019}{1}{Several companies occupied central positions within}. Table~6 is a worked re-identification scenario: a composite Alphabet profile assembled from three medicines apps (drug list, conditions, location) \pc{grundy2019}{9}{Risks to privacy owing to data aggregation within fourth party network}.

\cf{Limitations stated} Single time point, small purposive sample, Android only. Recipient characterisation relies on self-reported practices and the authors' interpretation, mitigated by duplicate extraction \pc{grundy2019}{9}{Our characterisation of the main}.

\cf{Use for this thesis} Mainly for \stagetag{H}: an evidence-based way to derive \emph{actors and access paths} (first/third/fourth parties) and \emph{aggregation capability} for a harm tree. Table~6 is a model of a concrete harm pathway narrative. For \stagetag{A}: triangulated purposive sampling (crawler plus endorsements) as an alternative to pure store rank. For \stagetag{D}: duplicate extraction of company documents into a structured form, and noting \emph{when} documents were captured relative to regulatory change. Traffic analysis is out of scope.
\end{papercard}
